\documentclass[12pt,a4paper]{article}

% =========================
% Packages
% =========================
\usepackage[margin=1in]{geometry}
\usepackage{setspace}
\usepackage{authblk}
\usepackage{graphicx}
\usepackage{booktabs}
\usepackage{multirow}
\usepackage{amsmath}
\usepackage{amssymb}
\usepackage[numbers,sort&compress]{natbib}
\usepackage[colorlinks=true,linkcolor=blue,citecolor=blue,urlcolor=blue]{hyperref}
\usepackage{indentfirst}

% =========================
% Basic formatting
% =========================
\onehalfspacing
\setcounter{secnumdepth}{3}
\setcounter{tocdepth}{3}

% =========================
% Outline placeholder control
% =========================
% Keep this TRUE while drafting the outline.
% Change to \showplaceholdersfalse after adding real body text.
\newif\ifshowplaceholders
\showplaceholderstrue

\newcommand{\outlineplaceholder}{%
\ifshowplaceholders
\par\noindent\textit{[Content to be developed.]}\par\medskip
\fi
}

% =========================
% Title information
% =========================
\title{\textbf{Embodied Intelligence for Next-Generation Medical Detection Technologies: Toward Closed-Loop Diagnostic Agents}}

\author[1]{Author One}
\author[1,2]{Author Two}
\author[1]{Corresponding Author\thanks{Corresponding author: \href{mailto:corresponding.author@email.com}{corresponding.author@email.com}}}

\affil[1]{Department or Institute Name, University or Hospital Name, City, Country}
\affil[2]{Department or Institute Name, University or Hospital Name, City, Country}

\date{}

\begin{document}

\maketitle

% =========================
% Abstract
% =========================
\begin{abstract}
% Medical detection technologies are undergoing a transition from centralized, episodic and passive testing toward distributed, continuous, adaptive and context-aware diagnostic systems. At the same time, embodied intelligence is shifting artificial intelligence from disembodied data interpretation toward agents that can perceive, reason, remember, interact with environments and initiate clinically meaningful actions. This Review examines the convergence of embodied intelligence and next-generation medical detection technologies. Rather than treating medical detection hardware and embodied artificial intelligence as separate domains, we propose an integrated framework in which sensing, sampling, measurement, interpretation, memory, action and feedback are coupled into closed-loop diagnostic systems. We discuss how embodied artificial intelligence can augment point-of-care testing, wearable and implantable biosensors, lab-on-chip platforms, automated laboratories, robotic sampling systems and multimodal diagnostic infrastructures. We then examine focused application landscapes in acute and hospital-based detection, decentralized and near-patient testing, home-based continuous monitoring and autonomous detection laboratories. Finally, we analyze validation, clinical translation, regulation, ethics and safety challenges, and identify three grand challenges for the next decade: self-optimizing diagnostic workflows, patient-specific embodied diagnostic models and closed-loop detection--intervention ecosystems. We argue that the central opportunity of embodied intelligence in medical detection is not merely to improve test interpretation, but to transform diagnostic technologies from passive measurement tools into adaptive diagnostic agents.
Abstract
\end{abstract}

\noindent\textbf{Keywords:} Embodied intelligence; Medical detection; Diagnostic agents; Wearable biosensors; Laboratory automation; Multimodal artificial intelligence

\tableofcontents
\newpage

% =========================
% Main text
% =========================

\section{Introduction: From Passive Testing to physical diagnostic agency}
\label{sec:introduction}

\subsection{The evolution of diagnostic paradigms and persistent gaps}
\label{subsec:evolution_medical_detection}

Diagnostic technologies underpin clinical decisions by transforming molecular, cellular, physiological and imaging signals into clinically interpretable evidence. Many diagnostic pathways are organized around discrete testing episodes: a patient or specimen is brought to a laboratory or specialist unit, measurements are obtained under controlled conditions, and the resulting report is interpreted within a wider care pathway. This model supports analytical rigour, standardization and quality control and remains indispensable for screening, diagnosis and therapeutic monitoring. Its output, however, is commonly a result associated with a particular patient or specimen, location and time.

As clinical demands have expanded towards earlier detection, rapid triage, longitudinal monitoring and decentralized care, limitations in this episodic model have become apparent and have helped drive successive technological responses. Temporal and spatial gaps in conventional testing have encouraged the development of miniaturized, point-of-care and sample-to-answer platforms that move selected assays closer to patients and shorten the interval between sampling and interpretation. Improvements in analytical performance have also drawn attention to vulnerabilities elsewhere in the testing pathway, including test selection, patient preparation, specimen or signal acquisition, transport, communication and follow-up. Programmable microfluidics and laboratory automation increasingly connect and standardize operations that were previously separated. Meanwhile, as diagnostic evidence has become distributed across devices, settings and time, the interpretive problem has broadened from analysing isolated results to integrating molecular measurements, physiological time series, medical images, electronic health records, patient-reported information and wearable or ambient data. Multimodal artificial intelligence has emerged in response to this heterogeneity, enabling complementary evidence to be combined and current observations to be related to prior patient states.

Collectively, these advances have made diagnostic workflows faster, more accessible and more information-rich by extending observation through time, standardizing multistep operations and integrating heterogeneous evidence. Yet these capabilities remain distributed across predefined tasks and disconnected stages. Software-based diagnostic agents reduce this fragmentation by maintaining context, reasoning over evolving evidence, planning multistep tasks and invoking digital tools with less continuous human oversight. Their agency, however, remains largely informational: they may identify the need for another measurement, but a clinician, technician or separately controlled instrument must still obtain it. Embodied diagnostic intelligence extends this architecture into the physical workflow by linking reasoning and planning to sensors and actuators. Decisions can thereby modify sampling, probe positioning or device operation, with their effects sensed and returned in real time. This perception--action loop extends automation from digital coordination to adaptive execution, allowing embodied diagnostic agents to regulate where, when and how evidence is acquired as biological states and diagnostic uncertainty evolve.


\subsection{The emergence of embodied diagnostic intelligence}
\label{subsec:emergence_embodied_medical_detection_intelligence}

Embodied intelligence has advanced through the convergence of sensorimotor learning, multimodal foundation models and predictive modelling. Early learned policies mapped observations directly to actions, while sequence-based and generative policies improved the precision and temporal coherence of execution. Vision--language--action models subsequently introduced semantic grounding, instruction following and broader task transfer, whereas video-prediction and world-model approaches enabled agents to anticipate how actions would change future observations. Recent architectures increasingly integrate understanding, future-state prediction and action generation within a single or tightly coupled system. These overlapping directions collectively mark a shift from reactive, task-specific control towards agents that can interpret goals, anticipate consequences and adapt execution through feedback.Against the acquisition and workflow gaps outlined above, this shift offers a natural route for making diagnostic sensing adaptive: physical acquisition can be adjusted according to signal quality, anatomical variation, diagnostic uncertainty and safety constraints, with the resulting evidence guiding subsequent actions.

Medical embodied intelligence has developed most visibly in surgical and interventional robotics. Autonomous laparoscopic anastomosis and image-guided needle steering have demonstrated perception-guided execution, replanning and recovery during interaction with deformable anatomy. Although primarily procedural, these systems establish capabilities directly relevant to diagnostic navigation, targeting and sampling. Their extension into medical detection is clearest in robotic ultrasound. Autonomous thyroid and carotid systems have integrated anatomical localization, probe and force control, image-quality assessment and disease-related interpretation within physical feedback loops. Related developments in robotic endoscopy, palpation and specimen acquisition suggest a wider transition from executing predefined medical actions towards actively regulating how diagnostic evidence is obtained. Most existing systems remain device- and task-specific, but collectively they provide the foundations for embodied agents that coordinate physical acquisition around evolving diagnostic needs.


\subsection{Scope and thesis of this Review}
\label{subsec:scope_thesis}

Building on the diagnostic gaps and technological trajectory outlined above, this Review adopts the complete diagnostic system, rather than any individual sensor, instrument or algorithm, as its primary unit of analysis. Its scope encompasses embodied intelligence across clinical care, diagnostic laboratories and everyday environments, while drawing selectively on surgical robotics, autonomous experimentation and assistive technologies where they provide capabilities transferable to medical detection. We first establish a conceptual framework for embodied medical detection and then examine the physical, interface and cognitive foundations required to implement it. Applications are subsequently organized by operational setting: clinical detection and intervention, laboratory diagnostics and assay development, and everyday monitoring and adaptive support. We further consider multilevel validation, clinical translation, lifecycle monitoring, ethics, safety and governance, before outlining three grand challenges for the field. Our central thesis is that the clinical value of embodied diagnostic intelligence will depend less on the autonomy of any single component than on the system-level integration of sensing, physical interaction, reasoning and workflow execution to generate reliable evidence, adapt safely to changing conditions and support accountable clinical decisions.


\section{Conceptual Framework: What Is Embodied Medical Detection?}
\label{sec:conceptual_framework}

\subsection{Defining embodied medical detection intelligence}
\label{subsec:defining_embodied_medical_detection}
Here, we define embodied medical detection as a system-level form of diagnostic intelligence in which sensing and interpretation are coupled to situated action within a patient, specimen, instrument or care environment. Three features are essential: the system must sense biologically or operationally relevant states; current evidence must influence an action that changes how further evidence is acquired, prepared, measured or verified; and the consequences of that action must inform what happens next. Such actions may include repositioning a probe, altering sampling or assay conditions, repeating a measurement or escalating an unresolved case. They may be executed by a device or mediated by a human, provided that they are causally integrated into the feedback process.

This definition distinguishes embodiment from physical proximity, automation and AI-enabled interpretation. A wearable sensor that samples on a fixed schedule, a model that interprets a completed dataset and a robot that follows a predetermined protocol can all support diagnosis without constituting embodied diagnostic intelligence. Embodiment arises only when these components form a feedback relation in which interpretation changes subsequent interaction and the resulting observations update the system. It is therefore a property of the integrated diagnostic system rather than of any sensor, robot, assay platform or algorithm in isolation.

Embodiment should also not be equated with unrestricted autonomy. Because diagnostic actions differ in risk and reversibility, authority may be distributed among automated control, operator confirmation and clinician oversight. A system may simply identify an inadequate signal and guide reacquisition, or it may coordinate multiple measurements and propose escalation when uncertainty remains. Across this continuum, the defining criterion is not independence from humans, but the capacity to adapt interaction with the environment in pursuit of diagnostically sufficient, reliable and safe evidence.

\subsection{Medical detection as a workflow}
\label{subsec:medical_detection_workflow}
Medical detection is a process through which a clinical objective is translated into evidence that can support a decision, rather than an isolated act of measurement. Across diagnostic settings, this process comprises a common set of interdependent functions: establishing the clinical objective and relevant context, acquiring evidence through sensing or sampling, preparing and measuring that evidence, assessing its quality and uncertainty, interpreting the resulting findings, and communicating a result or initiating action. These functions define a conceptual architecture rather than a fixed chronology. They may overlap, recur or be distributed across different people, instruments and settings, while safety, provenance and clinical context must be maintained throughout. The same functions are realized through different physical and informational operations across diagnostic modalities. Laboratory testing transforms biospecimens through collection, transport, preparation and assay. Medical imaging couples protocol and view selection with patient or probe positioning, signal acquisition and reconstruction, whereas physiological monitoring relates sensor placement and signal integrity to changes in biological state over time. Point-of-care, wearable and automated laboratory platforms may consolidate several functions within a single system, but they do not eliminate the dependencies among them. The meaning of a diagnostic result therefore remains inseparable from the state examined, the conditions under which evidence was acquired and transformed, the criteria used to establish its adequacy and the decision it is intended to inform.

Representing medical detection in this way provides a common basis for comparing otherwise heterogeneous technologies without reducing them to a particular device, assay or data modality. It directs attention to the interfaces through which biological material, signals, information and decision authority pass among patients, professionals and machines, as well as to the dependencies through which one function shapes the next. This functional architecture provides the scaffold for mapping perception, memory, reasoning and action onto medical detection in the following subsection.

\subsection{The embodied diagnostic loop}
\label{subsec}

Within this functional architecture, perception, memory, reasoning and action describe interdependent capacities distributed across the diagnostic workflow rather than discrete modules arranged in a fixed sequence. Perception encompasses the acquisition of signals from patients, specimens, instruments and care environments, together with the conditions that determine their quality. Memory retains prior observations, patient-specific baselines, acquisition histories and the consequences of earlier interventions, providing the context in which new evidence is interpreted. Reasoning relates this evolving evidence to the clinical objective, evaluates uncertainty and determines whether the available information is sufficient or whether another operation is required. Action may reposition a sensor or probe, collect or manipulate a specimen, adjust acquisition or assay conditions, select a complementary measurement or escalate the case. Crucially, the consequences of these actions are fed back through subsequent observations. These observations may confirm, refine or overturn the previous interpretation and determine whether the diagnostic objective has been met or further action is needed.

For this feedback to be diagnostically meaningful, the system must distinguish changes in the biological state under investigation from changes in the measurement process through which that state is observed. An altered signal may reflect genuine biological change, but may also arise from sensor displacement, specimen inadequacy, calibration drift or delayed data transmission. Repositioning the sensor, reacquiring the specimen or modifying the assay generates new evidence with which these possibilities can be resolved. Quality assessment therefore remains active throughout the diagnostic process, enabling successive observations and operations to reduce uncertainty and guide the workflow towards its clinical objective.


\section{Technical Foundations of Embodied Diagnostic Agents}
\label{sec:technical_enablers}

% Paragraph group 1:
% Physical and virtual embodiments

% Paragraph group 2:
% Sim-to-real translation, interfaces and connectivity

% Paragraph group 3:
% From teleoperation and predefined automation to adaptive agency

% Paragraph group 4:
% Capability continuum and transition to application domains

Embodied diagnostic agents depend on three coupled foundations: embodiment provides situated sensing and action, connectivity links observations, commands and outcomes across components, and adaptive agency uses feedback to adjust the diagnostic process. Their integration, rather than the capability of any component alone, determines whether a functional diagnostic loop can be sustained.

In the development of embodied diagnostic agents, virtual simulation provides a scalable environment for robot learning and evaluation before physical deployment. General-purpose simulators such as MuJoCo, Gazebo and Isaac Sim reproduce robot dynamics, sensing and physical interactions, while task-oriented platforms such as RoboTwin, ManiSkill and RLBench further provide standardized manipulation tasks, demonstrations, domain randomization and evaluation protocols for training and comparing imitation-learning, reinforcement-learning and vision--language--action policies. Medical-robotics frameworks such as SOFA, AMBF and ORBIT-Surgical extend these capabilities to deformable anatomy, tissue--instrument interactions and procedure-specific tasks. For diagnostic applications, simulation must additionally represent anatomical or specimen variability, contact-dependent signal quality and modality-specific sensing, while evaluation should consider diagnostic adequacy, physical safety, failure recovery and sim-to-real robustness rather than task completion alone. The resulting policies must ultimately be transferred to physical embodiments, including patient-facing sensors, point-of-care systems, sampling robots and laboratory manipulators, where performance is constrained by real anatomy, materials and clinical workflows. Virtual platforms therefore support scalable learning and predeployment evaluation, whereas physical embodiments determine whether diagnostic agency can be realized reliably in practice.

Connectivity couples virtual models, physical devices and clinical information systems through three technical mechanisms. Sim-to-real transfer and model recalibration address anatomical, sensor and instrument variation. Interoperable interfaces, including SMART on FHIR, preserve identity, timing, calibration, provenance and uncertainty as data move across components. Bidirectional control interfaces transmit requested operations, confirm their execution, return observed outcomes and fail safely when actions cannot be completed.

On this technical substrate, adaptive agency comprises two coupled functions: diagnostic decision-making and action execution. The diagnostic function updates the working diagnostic state, assesses evidence quality and sufficiency, and selects the next evidence-acquisition goal; the execution function realizes that goal through skill decomposition, motion or device planning, feedback control, physical actuation and verification. Both functions may be implemented by predefined rules, learned policies or hybrid systems. Thus, a fixed scanning or sampling sequence automates execution, whereas an observation-conditioned rule that repeats, redirects, escalates or terminates acquisition constitutes rule-based diagnostic decision-making. Model architectures distribute these functions differently. Modular systems may combine a pretrained vision--language model with an action expert; video or world-model approaches may predict action-conditioned futures to support action selection; and unified vision--language--action models may jointly learn observation--instruction--action mappings. Demonstration-based methods provide a complementary route for acquiring execution policies. These architectures do not by themselves determine autonomy: teleoperation, task autonomy and conditional or supervised autonomy instead describe how authority is allocated between the human and the system across diagnostic and executional functions. Nor does a model output constitute a completed action. Diagnostic updating may be event-driven, whereas motion, force and emergency-control loops require faster, bounded-latency feedback; VLA policies must therefore operate through local controllers, safety monitors, actuators and end effectors.

Closed-loop diagnostic agency arises when the observations and execution status resulting from one acquisition update the diagnostic state and determine whether evidence acquisition should continue, be repeated, redirected, escalated or stopped. Healthcare-agent architectures formalize such sequential coordination, while AMIE and MIRA demonstrate informational versions through adaptive history taking, test selection and iterative diagnostic reasoning in simulated dialogue or EHR environments. These systems remain disembodied unless their requests are coupled to patient-facing sensors or instruments; when a human performs the requested examination, the loop is physically grounded but human-mediated rather than autonomously embodied. Conversely, closed-loop insulin delivery and bidirectional neuromodulation demonstrate integrated sensing, control and actuation, but optimize therapeutic regulation rather than diagnostic information acquisition. Embodied diagnostic agency lies at the intersection: an explicit diagnostic objective guides situated sensing or sampling, the execution stack performs and verifies the operation within temporal and safety constraints, and the returned evidence changes the next decision. A VLM, video or world model, or action model therefore contributes to diagnostic agency only when its outputs govern evidence acquisition and subsequent decisions are revised using the consequences of action. Clinically consequential actions must additionally be constrained by clinical relevance, evidence sufficiency, expected information gain, risk, and explicit human authorization and override.

These foundations define a function-specific continuum of autonomy rather than a path towards unrestricted autonomy. Authority should be allocated according to the time sensitivity, uncertainty, reversibility and potential harm of each function: low-risk acquisition steps may be automated, whereas invasive sampling, clinical escalation and other consequential actions remain subject to clinical authorization. Clinical readiness must therefore be assessed at the level of the integrated sensing--decision--action--feedback loop rather than inferred from the isolated performance of a sensor, model or robotic task. Evaluation should examine sequential operation under changing observations, calibration drift, missing or delayed information, unsuccessful actions, interface failures and human intervention, including whether the system detects failure, degrades safely, recovers or transfers control. Although clinical evidence for such integrated embodied diagnostic loops remains limited, this closed-loop organization and its allocation of authority define the technical design space for the patient-facing, laboratory and everyday applications considered below.



\section{Application landscapes}
\label{sec:application_landscapes}

% Opening paragraph:
% Introduce the application landscape according to the environments in
% which diagnostic agents are embedded and where their
% perception--reasoning--action loops are closed.
%
% The following applications are organized into three interconnected
% settings: clinical care, diagnostic laboratories and everyday life.
% These settings should not be viewed as isolated device categories,
% but as successive and interacting components of a longitudinal
% diagnostic pathway.
The application landscape of embodied diagnostic agents is best distinguished by the environments in which their perception--reasoning--action loops are closed rather than by device form alone. We organize this landscape into three interconnected settings: clinical care, diagnostic laboratories and everyday life. Clinical environments centre the loop on a patient-specific question and the acquisition of evidence that can guide an immediate examination, test or intervention. Diagnostic laboratories relocate the loop to specimens and instruments, where analytical operations are selected and adapted according to sample state and intermediate results. Everyday systems extend observation beyond formal encounters and may determine when clinical assessment or laboratory testing should be initiated. These settings form successive and interacting components of a longitudinal diagnostic pathway rather than isolated device categories. A signal detected at home may trigger clinical evaluation; a patient-facing procedure may generate a specimen for centralized analysis; and a laboratory result may redirect bedside examination, treatment or subsequent monitoring. The interfaces between settings are therefore part of the diagnostic loop: clinically meaningful continuity depends on whether patient identity, specimen integrity, acquisition context, uncertainty and urgency remain connected as evidence and actions move across environments.


\subsection{Clinical detection and intervention}
\label{subsec:clinical_detection_intervention}

% Paragraph 1: Patient-facing acquisition and assessment
%
% Begin with a clinical question or an initial diagnostic hypothesis,
% which determines what evidence should be acquired from the patient.
%
% Discuss embodied systems that directly interact with patients through
% robotic or image-guided sampling, bedside examination, continuous
% physiological sensing and intra-procedural assessment.
%
% Focus on the clinical need for safe and reliable evidence acquisition
% under patient movement, anatomical variability, time pressure and
% uncertainty.
%
% Explain how these systems perceive patient anatomy, physiological state
% or signal quality; select or adjust an acquisition strategy; perform the
% measurement or sampling action; and verify whether diagnostically useful
% evidence has been obtained.
%
% Representative examples may include robotic blood collection,
% image-guided biopsy, AI-assisted point-of-care ultrasound,
% robotic auscultation, continuous inpatient monitoring and
% intra-procedural sensing.


% Paragraph 2: Distributed testing and diagnostic continuity
%
% Explain that patient-facing acquisition can lead to two diagnostic routes.
% In one route, bedside or point-of-care systems immediately process and
% interpret the acquired information. In the other, physical specimens are
% transferred to a centralized diagnostic laboratory.
%
% For point-of-care testing, discuss the integration of sample preparation,
% measurement, quality assessment and machine-learning-assisted
% interpretation close to the patient.
%
% For centralized testing, discuss patient--sample identification,
% labelling, sample-quality assessment, chain-of-custody tracking and
% autonomous specimen transport from wards, clinics or operating rooms
% to the laboratory.
%
% Emphasize that embodied diagnostic logistics should preserve the identity,
% integrity, context and clinical priority of the specimen, rather than
% merely automate physical transportation.
%
% Define the clinical--laboratory boundary:
% this paragraph ends at formal laboratory accessioning, whereas subsequent
% specimen processing and analytical operations are discussed in the
% laboratory subsection.


% Paragraph 3: Detection-guided decision, action and reassessment
%
% Describe how patient observations, bedside measurements, imaging,
% laboratory results and longitudinal records are integrated into a
% clinically meaningful representation of the patient's state.
%
% Discuss sequential diagnostic reasoning, including the selection of
% additional questions, examinations or tests when current evidence is
% incomplete or uncertain.
%
% Present the clinical closed loop as:
% clinical question
% $\rightarrow$ targeted evidence acquisition
% $\rightarrow$ multimodal interpretation
% $\rightarrow$ selection of the next test or clinical action
% $\rightarrow$ reassessment using newly acquired evidence.
%
% Distinguish between measurement actions, diagnostic workflow actions
% and therapeutic interventions. Emphasize that higher-risk clinical
% actions generally require explicit clinician supervision, clear
% escalation criteria and the ability to override or stop the system.
%
% Clarify that a standalone diagnostic model is not necessarily embodied.
% Embodied diagnostic agency emerges when perception and reasoning are
% connected to physical diagnostic tools, workflow operations or
% clinically consequential actions.
Clinical detection begins with a question or provisional diagnostic hypothesis that determines what evidence should be acquired from the patient. An embodied acquisition system must translate this objective into a feasible and safe operation despite patient movement, anatomical variability, limited time and uncertainty about whether the resulting signal or specimen will be diagnostically adequate. It may localize an anatomical target, assess physiological state or signal quality, select an acquisition trajectory or sensor configuration, execute or guide the measurement and then verify the evidence obtained. Automated venipuncture illustrates this structure: ultrasound imaging is used to identify a suitable peripheral vessel, after which a robotic mechanism guides the needle towards the lumen and confirms the outcome of the attempt. In AI-guided echocardiography, a deep-learning system evaluates the current view and provides real-time instructions for probe movement, enabling novice operators to acquire images suitable for limited diagnostic assessment. Here, embodiment is distributed between the software and the human operator who physically executes its recommendations. A fully autonomous thyroid-ultrasound system integrates body-region localization, force feedback and adaptive probe orientation within the robotic platform itself, providing a more complete acquisition loop in which observations of anatomy and contact conditions continually modify probe motion. The same functional pattern applies to image-guided biopsy, robotic auscultation and intra-procedural sensing: the relevant achievement is not merely automated motion or signal classification, but the reliable acquisition of evidence that answers the clinical question. Continuous inpatient wearables extend this process over time by repeatedly sampling physiological state and enabling deterioration risk to be updated as new observations arrive. Nevertheless, a wearable prediction model alone does not constitute an embodied diagnostic agent unless its output is connected to actions such as checking sensor placement, increasing measurement frequency, requesting confirmatory assessment or escalating the patient for clinical review.

Patient-facing acquisition subsequently gives rise to two diagnostic routes. In the first, information or specimens are processed near the patient, allowing bedside or point-of-care systems to combine sample preparation, measurement, quality assessment and algorithmic interpretation within a short local workflow. Microfluidic platforms have, for example, quantified leukocyte counts and neutrophil CD64 expression directly from small volumes of whole blood without manual sample processing. More recent systems integrate automated whole-blood preparation, molecular measurement and machine-learning-derived classifiers to estimate deterioration risk in patients with suspected sepsis. Machine learning can additionally support signal reconstruction, assay reading, quality control and interpretation across multiple point-of-care modalities. Such integration shortens the distance between acquisition and interpretation, but the loop closes locally only if the system can recognize an inadequate specimen or invalid measurement and initiate an appropriate response, such as repeating the assay, requesting another sample or escalating the result. In the second route, a physical specimen is transferred to a centralized laboratory. During this interval, the specimen becomes the material carrier of evidence about the patient, and preanalytical failures in identification, labelling, collection, handling or transport can invalidate otherwise accurate downstream analysis. Embodied diagnostic logistics must therefore do more than navigate between a ward and a laboratory: they must maintain the link between patient, order and specimen; monitor relevant transport conditions and delays; document handovers; and preserve clinical priority. Hospital logistics robots can combine autonomous navigation with authenticated pickup and receipt, status tracking, hospital-information-system integration and urgency-based task scheduling, illustrating how transport can become part of a traceable diagnostic workflow. The clinical side of this pathway ends at formal laboratory accessioning, when receipt, identity, condition and acceptance of the specimen are recorded. Subsequent specimen preparation and analytical operations belong to the laboratory environment discussed below.

Evidence acquired through examination, monitoring, imaging and laboratory testing must ultimately be assembled into a clinically meaningful representation of the patient's evolving state. Because the initial evidence is often incomplete or ambiguous, diagnostic reasoning is sequential: each result can change the differential diagnosis and determine which question, examination or test would be most informative next. Conversational diagnostic systems demonstrate active history taking in which successive questions depend on previous answers, while differential-diagnosis models and multimodal oncology agents illustrate how clinical records, images, specialised analytical tools and external knowledge can be coordinated during interpretation. These systems primarily establish agency at the informational level; they become clinically embodied only when their selections can alter subsequent acquisition or care. The corresponding clinical loop can be expressed as

The actions within this loop differ in consequence. Measurement actions adjust a sensor, reacquire an image or repeat a reversible test. Diagnostic-workflow actions request another examination, prioritize a specimen, issue an alert or escalate the case to a clinician. Therapeutic interventions alter the patient's biological state and therefore carry a different level of risk. The deployment of the TREWS sepsis warning system illustrates a clinician-mediated workflow loop in which model-generated alerts were reviewed by providers and connected to subsequent clinical decisions, rather than executed autonomously. Closed-loop insulin delivery provides a mature example of direct sensing, algorithmic control, physiological intervention and feedback, but it is a therapeutic control system rather than a diagnostic agent. As actions become more invasive, less reversible or more clinically consequential, explicit clinician authorization, defined escalation criteria and the ability to interrupt or override the system become increasingly important. A standalone model that predicts a diagnosis from already collected data is therefore not necessarily embodied. Embodied diagnostic agency arises when perception and reasoning are connected to physical diagnostic tools, workflow operations or clinically consequential actions, and when the results of those actions are observed and used to determine what should happen next.


\subsection{Laboratory diagnostics and assay development}
\label{subsec:laboratory_diagnostics}

% Paragraph 1: From manual manipulation to programmable execution
% Repetitive physical operations, occupational burden, robotic liquid
% handling and task-level automation.

% Paragraph 2: From fixed workflows to adaptive laboratory operation
% Integration of instruments, perception of sample and equipment states,
% adaptation to variable conditions, error recovery and dynamic replanning.

% Paragraph 3: From automated experimentation to closed-loop scientific agency
% Scientific objectives, experiment selection, robotic execution,
% measurement, reasoning, model updating and selection of the next experiment.
After formal accessioning, the centre of embodiment shifts from the patient to the specimen and the instruments that transform it into diagnostic evidence. Many preanalytical and analytical workflows nevertheless remain dependent on repetitive manual operations, including tube sorting, centrifugation, decapping, aliquoting, reagent transfer, plate handling and instrument loading. Programmable automation converts these operations into reproducible machine-executable protocols while using barcodes and laboratory information systems to preserve the association among the patient, test order, specimen and derived aliquots. An early evaluation of an automated preanalytical workstation, for example, demonstrated integrated specimen sorting, centrifugation, decapping, labelling, aliquoting and rack loading, with reduced manual staffing requirements and fewer specimen-processing errors. Total laboratory automation extends this principle by linking preanalytical modules, analysers and postanalytical storage or retrieval through conveyors, robotic workcells and middleware. In research and assay-development laboratories, automated liquid handlers similarly enable larger and more precisely timed experiments across microlitre and sub-microlitre volume ranges. Their performance, however, remains dependent on liquid properties, transfer volume, labware geometry, calibration and the compatibility of the surrounding protocol. Fixed automation is therefore most effective for standardized and repetitive operations; unusual specimens, non-standard containers, method changes, equipment maintenance and workflow exceptions continue to require human expertise. At this stage, the principal gain is executional scalability rather than diagnostic or scientific autonomy: the system can reliably perform what a validated protocol specifies, but it does not necessarily determine whether the protocol remains appropriate for the observed specimen or result.

Programmable execution becomes adaptive laboratory operation when information about the specimen, instrument and workflow state is used to condition what happens next. Relevant observations may include specimen identity and volume, integrity indicators, analytical flags, quality-control status, device availability, transfer success and the plausibility or consistency of intermediate results. Autoverification provides an established rule-based example: predefined criteria allow routine results that satisfy analytical and consistency checks to be released automatically, while discordant, flagged or otherwise exceptional cases are routed for specialist review. This illustrates that laboratory decision-making need not depend on a learning model; observation-conditioned rules that repeat a measurement, hold a result, redirect a specimen or request human review already create a bounded operational feedback loop. More advanced laboratory automation seeks to embed real-time quality control, specimen-quality assurance and result verification throughout the workflow rather than treating them as isolated post hoc checks. Physical adaptability additionally requires orchestration across heterogeneous instruments and automation islands. Hierarchical workflow-management systems can coordinate automated stations, mobile-robot transport and remaining manual operations, whereas plug-and-play architectures use standardized device descriptions, interfaces and robot-relevant spatial information to reduce the reprogramming required when instruments or layouts change. Perception can further close the execution loop. The LIRA module, for example, uses vision-based localization and a vision--language model to inspect rack or plate placement, classify manipulation errors and invoke predefined correction or human intervention. Such capabilities move laboratory robots beyond open-loop execution, but they should not be confused with unrestricted autonomy: current recovery strategies remain task-specific, and end-to-end integration among inspection, workflow scheduling and experimental replanning is still incomplete. Robust adaptive operation therefore depends on standardized interfaces, traceable state representations, bounded response times, exception handling and safe degradation when observations are ambiguous or an operation cannot be verified.

Assay development introduces a scientifically distinct closed loop. Here, the objective is not simply to process a patient specimen according to a validated standard operating procedure, but to determine which experiment should be performed next to improve a reagent, protocol or analytical design. Automated experimentation forms an important intermediate stage. A robotic liquid-handling platform for lateral-flow assay development, for example, performed experiments directly in the LFA format, screening discrete variables such as antibody pairs and continuous variables such as reagent concentrations for malaria, SARS-CoV-2 and tuberculosis assays. By integrating liquid delivery, imaging and signal quantification, the platform increased experimental scale and reproducibility while reducing hands-on work. Its experimental space and optimization procedure, however, were principally specified by researchers; high-throughput execution alone therefore does not constitute a self-driving laboratory. Scientific agency emerges when measurements are automatically interpreted, used to update a model of the experimental space and converted into the next experimental choice. The resulting loop links a scientific objective to candidate selection, robotic preparation and testing, quality-controlled measurement, model updating and selection of the subsequent experiment. Chemical systems have demonstrated this principle through Bayesian selection of experiments, language-model-assisted planning and the integration of multiple analytical instruments. Closer to biomedical assay development, the SAMPLE platform coupled protein-sequence design with automated DNA assembly, protein expression and biochemical measurement, using each experimental cycle to update a sequence--function model and select new protein variants. Similar closed loops could eventually optimize diagnostic recognition elements, reagent combinations, assay conditions and decision thresholds. Nevertheless, optimization of analytical signal or molecular function does not by itself establish clinical validity. Candidate assays must still be evaluated using appropriate positive, negative and process controls; clinically representative matrices and independent specimens; interference and cross-reactivity studies; lot, instrument and site variation; prespecified performance criteria; and independent analytical and clinical validation. The assay-development loop and the patient-level diagnostic loop should therefore remain conceptually distinct: a self-driving laboratory generates and evaluates candidate assays, whereas only a subsequently frozen, validated and governed assay can enter routine clinical, point-of-care or everyday diagnostic use.


\subsection{Everyday embodied monitoring and adaptive support}
\label{subsec:everyday_monitoring}

% Intelligent body-integrated monitoring
% Wearable, implantable and ingestible systems

% Ambient and home-embedded intelligence
% Smart-home sensing and connected home diagnostic devices

% Interactive care and functional assistance
% Companion robots, rehabilitation robots and intelligent assistive devices
Beyond episodic encounters and centralized laboratories, everyday monitoring relocates evidence acquisition to the body and the environments in which health changes unfold. The aim is not merely to collect more measurements, but to learn an individual baseline, detect meaningful departures from it and initiate proportionate verification or support. Wearable systems provide the most accessible implementation, combining a body interface with sensing, power, communication, signal processing and a decision unit; failures in any one component can invalidate the end-to-end inference. Consumer smartwatches have, for example, detected presymptomatic physiological deviations relative to personal baselines in COVID-19, while multimodal wearable data have supported personalized prediction of selected clinical laboratory measurements. Biochemical wearables extend observation beyond motion and vital signs: an integrated electrochemical platform combined induced sweat, microfluidic handling, calibration and wireless analysis to monitor trace metabolites and nutrients during exercise and at rest. Implantable and ingestible devices gain access to signals that are difficult to sample reliably at the skin. A microfluidic intraocular-pressure implant enabled smartphone-based home readings across fluctuations missed by occasional clinic measurements, and active-reset affinity sensors overcame saturation to track inflammatory proteins dynamically in living animals. Ingestible capsules have measured intestinal gases in a human pilot, self-powered glucose sensing in the porcine small intestine and, in a first-in-human study, respiratory and cardiac activity from within the gastrointestinal tract. These studies expand the spatial, temporal and molecular reach of observation, but their translational stages differ sharply. Motion artefact, biofluid variability, sensor fouling, calibration drift, data loss, power and telemetry constraints, foreign-body effects and uncertain links between proxy signals and clinical states must be evaluated longitudinally in the intended population. Most current platforms also remain one-way monitors: an alert or risk score becomes a diagnostic loop only when it triggers an appropriate repeat measurement, confirmatory examination or escalation, and the resulting evidence updates what happens next.

Ambient intelligence moves the sensing interface from the body into the home. Passive infrared detectors, door contacts and mattress sensors can generate a ``zero-interaction digital exhaust'' describing mobility, sleep, activity patterns and functional change without requiring the user to remember, wear or charge a device. Contactless radio-frequency sensing similarly enabled a model to infer Parkinson's disease and estimate severity from nocturnal breathing, including measurements acquired in the home. Connected fixtures can turn routine activities into sampling opportunities. A mountable smart-toilet prototype integrated user identification, urinalysis, uroflowmetry and stool-image analysis, illustrating how excreta could be monitored without a separate testing episode; however, passive collection raises unusually persistent questions about identity, consent, data ownership and access, especially in shared households. More broadly, home monitoring may support ageing in place and chronic-condition management, but the evidence base remains heterogeneous and digital access, maintenance and workflow integration can determine who benefits. The home is therefore not a controlled laboratory. Visitors, pets, co-residents, altered routines, moved devices, connectivity loss and gradual sensor drift can all change the meaning or completeness of a signal. Evaluation should consequently test person-level longitudinal calibration, missingness, false-alert burden and performance across households and social contexts, as well as whether a detected change leads to timely and useful action. In many near-term systems, the operational loop will remain human-mediated: ambient sensing identifies a deviation, a patient, caregiver or clinician verifies its context, and only then is monitoring intensified or care escalated.

Adaptive support begins when observations of the person change a physical or socially situated response. In rehabilitation, wearable sensing and machine learning can estimate motor impairment, movement quality and recovery outside periodic clinic assessments, creating a state representation with which therapy can be individualized. Electromechanical and robot-assisted arm training after stroke has shown modest improvements in activities of daily living, arm function and strength across randomized trials, although interventions, populations and training dose were heterogeneous and many devices delivered repetitive, predefined exercise rather than online adaptation. Human-in-the-loop optimization demonstrates a tighter embodied cycle: assistance parameters are varied, the user's metabolic or kinematic response is measured, and the controller selects a new policy for the same person. This principle has been extended from tethered laboratory systems to portable ankle exoskeletons that use wearable kinematics to personalize assistance during variable-speed outdoor walking. The engineering result is important, but optimization in small studies of controlled walking should not be equated with durable rehabilitation benefit, fall prevention or safe adaptation across impairment, fatigue and terrain. Socially assistive and companion robots occupy a different action space, using speech, movement, reminders and structured interaction to support engagement or well-being. Meta-analyses in older adults and people living with dementia suggest feasibility and possible benefits for selected emotional or neuropsychiatric outcomes, but effects on cognition, agitation and quality of life are inconsistent and high-quality randomized evidence remains limited. A recent randomized external pilot in surgical wards likewise demonstrated the feasibility of deploying a predefined-dialog robot, but found limited overall effects on engagement and perceived care quality. Across these devices, adaptation should therefore be judged against the intended functional or care outcome, not interaction frequency alone. Integrated evaluation must include state-estimation error, controller latency, physical safety, comfort and fatigue, adherence, human override, caregiver workload, privacy and inequitable fit. Everyday embodied support is thus a function-specific continuum: monitoring may remain passive, low-risk assistance may adapt locally, and consequential clinical or physical changes may require confirmation and authorization. Its defining advance is not continuous sensing or robotics in isolation, but a bounded sensing--decision--action--feedback loop that improves everyday function or care while preserving the user's agency and a safe path to human intervention.


\section{Evaluation, Validation and Clinical Translation}
\label{sec:validation_clinical_translation}
Clinical translation requires a continuous chain of evidence from valid measurement to reliable physical operation, effective human--system collaboration and sustained real-world performance. Because failures can arise at any interface among sensing, inference, action and human oversight, evidence for an isolated sensor, model or robotic task cannot by itself establish the safety or clinical value of the integrated system.


\subsection{Multilevel validation and benchmarking}
\label{subsec:multilevel_validation_benchmarking}
The second level evaluates embodied capability within a defined operating envelope. Benchmarks should reproduce the relevant anatomy, specimens, users, instruments and environmental disturbances, and compare the system with current practice, human or teleoperated operation and appropriate non-learning baselines. In addition to task success and completion time, evaluation should report retries, intervention and abstention rates, sensing or manipulation failures, force or exposure limits, specimen adequacy, contamination or damage, latency, severe-failure frequency and recovery. Requirements should become more stringent as autonomy and potential consequences increase: a system that merely recommends a sampling location requires different evidence from one that autonomously contacts tissue, obtains a specimen or escalates care. The staged development, comparative evaluation and long-term monitoring proposed by the IDEAL robotics framework provide a useful translation pathway beyond surgical applicationsd.

The third level evaluates the complete sensing--decision--action--feedback loop. Prospective tests should determine whether sequential use improves diagnostic adequacy, time to an appropriate decision, patient-relevant outcomes or resource use relative to the current workflow without increasing harm or burden. Evaluation should include changing observations, missing or delayed inputs, calibration drift, unsuccessful actions, human interruption and safe transfer of control. External validation should be aligned with the intended claim and distinguish temporal, geographical and domain generalizability \cite{deHond2023ClinicalAlgorithmValidation}. Credible benchmarks should therefore be clinically motivated, independently partitioned at the patient, specimen, operator, device and site levels where appropriate, and accompanied by reproducible protocols, meaningful baselines and a failure taxonomy rather than optimized solely for a leaderboard metric \cite{Mincu2022HealthcareBenchmarks}. Claims must remain bounded by the highest level reached: benchtop or simulated success supports technical feasibility, whereas clinical readiness requires evidence from the integrated system in its intended workflow.


\subsection{Human-centred clinical evaluation}
\label{subsec:human_centred_clinical_evaluation}
Human-centred evaluation treats the patient, clinician, technician or caregiver as part of the clinical system rather than as an external observer. Development should first identify the intended users, environments, critical tasks, foreseeable use errors, accessibility requirements and allocation of authority. Formative studies with representative stakeholders can refine workflow and interface design, followed by summative human-factors validation under realistic conditions. Early live evaluation should use a sufficiently stable system version and document user training, case selection, workflow integration, system outputs, human responses, overrides, technical failures and adaptations made during use, consistent with DECIDE-AI.

The principal comparator is the performance of the human--AI--robot team against the existing workflow, not the isolated model against an expert label. Relevant outcomes include diagnostic yield and error, task completion, time to appropriate action, unnecessary repetition or escalation, intervention and override quality, workload, situation awareness, usability, trust calibration, physical comfort, patient experience and staff or caregiver burden. Qualitative observation is valuable for identifying mode confusion, workarounds, communication failures and reasons for non-use that aggregate accuracy measures cannot reveal. Human--AI collaboration may reduce reading time or case volume in particular imaging workflows, but such benefits depend on how AI is positioned within the workflow and cannot be inferred from stand-alone accuracy. Analyses should also examine differential access, failure and benefit across relevant demographic, clinical, functional and professional subgroups.

When moving to prospective comparative or randomized studies, protocols should prespecify the algorithm and hardware versions, required user skills, input and output handling, human--system interaction, error analysis, outcomes and adverse-event procedures. SPIRIT-AI and CONSORT-AI provide corresponding guidance for trial protocols and reports. However, reporting checklists do not replace appropriate controls, statistical design or patient-relevant endpoints. For systems performing physical or consequential actions, evaluation must additionally establish when confirmation is required, how the system can be interrupted, who remains accountable and whether control can be transferred without creating a new hazard. Stakeholder participation, human autonomy, transparency, fairness and responsibility should therefore be incorporated throughout study design and governance rather than assessed only after technical development.


\subsection{Deployment and lifecycle monitoring}
\label{subsec:deployment_lifecycle_monitoring}
Deployment should be treated as a monitored clinical intervention rather than the endpoint of validation. A staged pathway can move from safety and technical feasibility to efficacy under controlled conditions, effectiveness relative to existing practice, and post-deployment monitoring. Before routine use, each site should verify local data and workflow compatibility, hardware and software integration, user training, access control, maintenance, escalation procedures and rollback or suspension mechanisms. The operational design domain, system version, intended users, excluded conditions and allocation of responsibility should remain explicit. These practices connect the fairness, universality, traceability, usability, robustness and explainability principles of FUTURE-AI to concrete lifecycle controls.

Monitoring should combine input, process, output and outcome indicators. Version-tagged records may include signal or specimen quality, missingness, population and case mix, calibration, subgroup performance, alerts, abstentions, human interventions, physical task failures, recovery attempts, adverse events, uptime and workload. For embodied systems, monitoring must additionally cover sensor drift, actuator or instrument wear, consumables, cleaning, environmental changes and failures at hardware--software or device--workflow interfaces. Medical algorithmic audits can connect detected errors to the clinical task, affected subgroups, contributing system components and corrective actions. Distribution-shift detectors may identify changes in image quality, population composition or comorbidity patterns, but a detected shift is an early warning rather than proof of degraded clinical performance; it should trigger outcome-linked investigation and root-cause analysis.

Every monitored signal should therefore have a predefined response, which may include recalibration, intensified review, restriction of the operating domain, retraining, rollback, user notification or temporary suspension. Clinical evaluation remains an ongoing process because real-world evidence may reveal performance that is either better or worse than predeployment estimates. Good machine-learning practice similarly emphasizes representative data, independent testing, human--AI team performance, traceability and monitoring across the total product lifecycle. Model updates must not occur as uncontrolled learning during clinical use. Where predetermined change-control mechanisms apply, the permitted modifications, data and retraining procedures, verification and validation methods, acceptance criteria and impact assessment should be specified prospectively. Lifecycle governance thus closes the translational loop: field evidence is converted into auditable corrective action while system evolution remains bounded by its validated intended use and applicable regulatory oversight.


\section{Ethics, Safety and Governance}
\label{sec:ethics_safety_governance}

\subsection{Privacy and autonomy in continuous detection}
\label{subsec:privacy_autonomy}
\outlineplaceholder

\subsection{Safety and accountability in diagnostic action}
\label{subsec:safety_accountability}
\outlineplaceholder

\subsection{Equity and resilience of connected systems}
\label{subsec:equity_resilience}
\outlineplaceholder


\section{Outlook: Three Grand Challenges}
\label{sec:future_outlook}

\subsection{Self-optimizing diagnostic workflows}
\label{subsec:self_optimizing_diagnostic_workflows}

\subsection{Patient-specific diagnostic models}
\label{subsec:patient_specific_embodied_models}

\subsection{Detection--intervention ecosystems}
\label{subsec:closed_loop_detection_intervention}


\section{Conclusion}
\label{sec:conclusion}
\outlineplaceholder

% =========================
% Declarations
% =========================

\section*{Author Contributions}
\outlineplaceholder

\section*{Declaration of Competing Interest}
\outlineplaceholder

\section*{Acknowledgements}
\outlineplaceholder

\section*{Funding}
\outlineplaceholder

\section*{Data Availability}
\outlineplaceholder

% =========================
% References
% =========================
% When you have a references.bib file and citations in the text, uncomment:
%
% \bibliographystyle{unsrtnat}
% \bibliography{references}

\end{document}
